\chapter{Cosa prevediamo}
\label{cap:previsioni}

Le previsioni si scrivono e si salvano su Git \textbf{prima} degli esperimenti: la data del salvataggio dimostra che non le abbiamo aggiustate dopo aver visto i risultati.

\section{Come si calcola un tempo di fermo}

\[
  \text{fermo} \approx \text{tempo per accorgersene} + \text{tempo per riavviare} + \text{tempo per tornare pronto}
\]

\begin{itemize}
  \item Se un programma \textbf{si spegne}, Kubernetes se ne accorge subito.
  \item Se \textbf{si congela}, se ne accorge solo al controllo successivo, quindi conta quanto spesso controlla.
  \item Con la \textbf{coda} (V3, V5) va aggiunto il tempo per smaltire gli ordini accumulati.
  \item Con il \textbf{database di riserva} (V4) il riavvio è sostituito dal passaggio alla riserva.
\end{itemize}

\section{Dalla misura alla disponibilità}

Ipotizziamo quanto spesso capita ogni guasto (per esempio, un'API che si spegne una volta al giorno) e moltiplichiamo per il fermo misurato. Sommando tutti i guasti otteniamo i minuti di fermo al mese, da confrontare con il limite di 43 minuti (D0017).

\section{Previsioni}

\begin{table}[htbp]
\centering
\caption{Previsioni da completare prima degli esperimenti.}
\label{tab:previsioni}
\begin{tabular}{@{}llcc@{}}
\toprule
\textbf{Guasto} & \textbf{Versione} & \textbf{Fermo previsto} & \textbf{Ordini falliti previsti} \\
\midrule
F1 & V1 & \todo{} & \todo{} \\
F1 & V2 & \todo{} & \todo{} \\
F4 & V1, V2 & \todo{} & \todo{} \\
F4 & V3 & \todo{} & \todo{} \\
F4 & V4 & \todo{} & \todo{} \\
F7 & V2 & \todo{} & \todo{} \\
\bottomrule
\end{tabular}
\end{table}
